\documentclass{oci}
\usepackage[utf8]{inputenc}
\usepackage{lipsum}

\title{Baldosas}

\begin{document}
\begin{problemDescription}
Varideri es el nuevo propietario de uno de los más famosos restaurantes de Portugal por sus hermosos diseños de azulejos.
En su primera inspección, notó que la calidad de la comida era notablemente inferior a la de otros restaurantes por la zona, por lo que decidió indagar en el asunto.

\setlength{\tabcolsep}{0.5pt}
\begin{figure}[h]
  \begin{center}
    {
      \renewcommand{\arraystretch}{0.25}
      \begin{tabular}{cccc}
        \includegraphics[scale=0.3]{a.eps} &
        \includegraphics[scale=0.3]{b.eps} &
        \includegraphics[scale=0.3]{a.eps} \\
        \includegraphics[scale=0.3]{b.eps} &
        \includegraphics[scale=0.3]{a.eps} &
        \includegraphics[scale=0.3]{b.eps} \\
        \includegraphics[scale=0.3]{a.eps} &
        \includegraphics[scale=0.3]{b.eps} &
        \includegraphics[scale=0.3]{a.eps}
      \end{tabular}
    }
  \end{center}
\end{figure}

En la inspección, notó que el embaldosado de la cocina está formado por dos tipos de azulejos, $0$ y $1$.
El diseño de estas baldosas está pensado para producir un patrón de cuadros alternados, es decir, cada baldosa de tipo $0$ tiene una baldosa de tipo $1$ junto a cada uno de sus lados y viceversa.
Sin embargo, los patrones en la cocina no respetaban esa regla, irritando así al chef y afectando la calidad de su comida.
Arriba puedes ver un embaldosado correcto, mientras que abajo aparece uno de los que están en la cocina.

\begin{figure}[h]
  \begin{center}
    {
      \renewcommand{\arraystretch}{0.25}
      \begin{tabular}{cccc}
        \includegraphics[scale=0.3]{a.eps} &
        \includegraphics[scale=0.3]{a.eps} &
        \includegraphics[scale=0.3]{b.eps} \\
        \includegraphics[scale=0.3]{b.eps} &
        \includegraphics[scale=0.3]{b.eps} &
        \includegraphics[scale=0.3]{a.eps} \\
        \includegraphics[scale=0.3]{a.eps} &
        \includegraphics[scale=0.3]{a.eps} &
        \includegraphics[scale=0.3]{b.eps}
      \end{tabular}
    }
  \end{center}
\end{figure}

Por esto, ha decidido reparar estos patrones para calmar el espíritu riguroso del chef y restaurar la calidad de sus platos.
Sin embargo, el presupuesto que Vardieri tiene disponible es limitado y necesita saber cuál es la forma más económica de realizar esta tarea.
Para cambiar una baldosa de tipo $0$ por una de tipo $1$ necesita desembolsar $P$, y para cambiar una de tipo $1$ por una de tipo $0$ necesita desembolsar $Q$.

¿Puedes ayudar a Varideri a encontrar la solución más económica para restaurar correctamente los azulejos?

\end{problemDescription}

\begin{inputDescription}
La entrada tiene varias líneas.

La primera línea contiene 4 enteros, $M$, $N$, $P$ y $Q$ ($1 \le M, N \le 10^3$, $1 \le P, Q \le 10^2$) que corresponden a las dimensiones $M \times N$ de un segmento rectangular de baldosas, y a los precios de cambiar las baldosas.

Las siguientes $M$ líneas contienen $N$ números, cuyos valores pueden ser $0$ ó $1$.
El $i$-ésimo número de la $j$-ésima que indican qué tipo de baldosa está en la posición (i, j) del embaldosado.
\end{inputDescription}

\begin{outputDescription}
La salida debe contener una línea con un único entero, el costo mínimo necesario para cambiar el embaldosado a un patrón de tablero de ajedrez.
\end{outputDescription}

\begin{scoreDescription}
  \subtask{20}
  Descripción Subtarea1
  \subtask{30}
  Descripción Subtarea2
  \subtask{50}
  Descripción Subtarea3
\end{scoreDescription}

\begin{sampleDescription}
\sampleIO{sample-1}
\sampleIO{sample-2}
\end{sampleDescription}

\end{document}
